\section{Уверенность рекомендаций и калибровка}

\subsection{Роль confidence-слоя в рекомендательной системе}

После перехода к рекомендательной постановке возникает вопрос не только о том, \textbf{какие} культуры попадают в
shortlist, но и о том, \textbf{насколько надёжно} модель формирует этот список. В прикладном сценарии это особенно
важно: два наблюдения могут иметь одинаковый формат выхода в виде Top-k кандидатов, но существенно различаться по
степени уверенности модельного вывода. Поэтому в работе поверх recommendation layer вводится дополнительный
confidence-слой. Его задача состоит в том, чтобы разделять случаи на более устойчивые и более спорные и тем самым делать
использование рекомендаций более интерпретируемым.

Содержательно confidence-слой не заменяет модель и не формирует новое ранжирование кандидатов. Он использует уже
рассчитанные вероятности CatBoost и переводит их в более прикладной формат. В результате пользователь получает не только
shortlist культур, но и дополнительный сигнал о том, к какому уровню надёжности относится конкретная рекомендация. Такая
логика особенно полезна в задачах поддержки принятия решений, где важно различать ситуации, в которых модель действует
на устойчивом статистическом основании, и ситуации, в которых результат следует интерпретировать осторожнее.

В работе используются три confidence-группы:

\begin{itemize}
\tightlist
\item
  \texttt{high\_confidence};
\item
  \texttt{medium\_confidence};
\item
  \texttt{low\_confidence}.
\end{itemize}

Пороговые правила формирования confidence-групп приведены в таблице~\ref{tab:confidence-thresholds}.

\begin{center}
\small
\begin{longtable}{p{0.24\textwidth} p{0.34\textwidth} p{0.34\textwidth}}
\caption{Пороги формирования confidence-групп}
\label{tab:confidence-thresholds}\\
\toprule
Группа & \texttt{top1\_proba} & \texttt{margin\_top1\_top2} \\
\midrule
\endfirsthead

\toprule
Группа & \texttt{top1\_proba} & \texttt{margin\_top1\_top2} \\
\midrule
\endhead

\bottomrule
\endlastfoot

\texttt{high\_confidence} & \texttt{top1\_proba >= 0.70} & \texttt{margin\_top1\_top2 >= 0.25} \\
\texttt{medium\_confidence} & \texttt{top1\_proba >= 0.50} & \texttt{0.25 >= margin\_top1\_top2 >= 0.12} \\
\texttt{low\_confidence} & \texttt{top1\_proba < 0.50} & \texttt{margin\_top1\_top2 <= 0.12} \\

\end{longtable}
\footnotesize
Источник правил: \path{notebooks/07_recommendation_confidence.ipynb} (правила high/medium/low confidence).
\normalsize
\end{center}

Тем самым итоговый вывод системы можно описать не только как ранжированный список кандидатов, но и как комбинацию двух
компонентов:

\begin{center}\(
\text{shortlist кандидатов} + \text{confidence signal}
\)\end{center}

Именно такое представление делает рекомендательный результат ближе к реальному сценарию использования.

\subsection{Результаты по confidence-группам}

Для test-выборки confidence-слой даёт следующие результаты:

\begin{center}
\small
\begin{longtable}{p{0.32\textwidth} r r r}
\caption{Качество рекомендаций по группам уверенности на тестовой выборке}
\label{tab:confidence-groups}\\
\toprule
Группа & Число наблюдений & \texttt{Accuracy@1} & \texttt{Hit@3} \\
\midrule
\endfirsthead

\toprule
Группа & Число наблюдений & \texttt{Accuracy@1} & \texttt{Hit@3} \\
\midrule
\endhead

\bottomrule
\endlastfoot

\texttt{high\_confidence} & 3 042 231 & 0.84147 & 0.97713 \\
\texttt{medium\_confidence} & 1 691 260 & 0.60931 & 0.94252 \\
\texttt{low\_confidence} & 1 043 523 & 0.42902 & 0.88987 \\

\end{longtable}
\end{center}


Как видно из таблицы, группы действительно различаются по качеству рекомендаций. В зоне \texttt{high\_confidence} модель
показывает наиболее устойчивый результат: top-1 точность превышает 0.84, а правильная культура попадает в Top-3 почти в
98\% случаев. В группе \texttt{medium\_confidence} показатели снижаются, однако остаются высокими для
shortlist-постановки. В группе \texttt{low\_confidence} top-1 точность уже заметно ниже, но даже здесь значение
\texttt{Hit@3} остаётся достаточно высоким. Это означает, что в сложных случаях модель часто ошибается в первой позиции,
но всё ещё включает фактическую следующую культуру в число нескольких наиболее вероятных кандидатов.

С практической точки зрения это разделение имеет понятную интерпретацию. В случаях высокой уверенности результат можно
рассматривать как сравнительно устойчивую рекомендацию. В случаях средней уверенности пользователь, скорее, получает
shortlist, внутри которого целесообразно сравнивать несколько вариантов. В случаях низкой уверенности модельный вывод не
следует скрывать, но его правильнее подавать как спорный случай, требующий более осторожной интерпретации и, возможно,
дополнительной экспертной проверки. Таким образом, confidence-слой делает рекомендательную систему более информативной,
поскольку показывает не только направление рекомендации, но и характер её надёжности.

\subsection{Интерпретация confidence-слоя}

Полученные результаты показывают, что confidence buckets не являются декоративным дополнением, а действительно отражают
различия в качестве модельного вывода. Это важный момент для общей логики работы. Если бы качество в группах
\texttt{high}, \texttt{medium} и \texttt{low} различалось слабо, confidence-сигнал не имел бы практического смысла.
Однако в рассматриваемой системе наблюдается явная стратификация по надёжности, что подтверждает полезность выделения
confidence-слоя как самостоятельного уровня архитектуры.

Особенно важно, что даже в группе \texttt{low\_confidence} модель сохраняет полезность именно как рекомендательная
система. Низкая уверенность в данном случае не означает, что результат полностью бесполезен; она означает, что top-1
интерпретация недостаточна и что корректнее рассматривать несколько альтернатив одновременно. Это хорошо согласуется с
общей постановкой работы: система ориентирована не на выдачу единственного ответа, а на формирование shortlist вероятных
вариантов с дополнительным контекстом для их анализа.

\subsection{Калибровка вероятностей}

Наличие confidence-слоя поднимает ещё один важный вопрос: насколько вероятности, выдаваемые моделью, действительно
согласуются с фактической частотой правильных ответов. Для этого в работе дополнительно рассматривается калибровка
вероятностей. Если модель хорошо откалибрована, то её confidence-сигналы можно интерпретировать не только как внутренний
score ранжирования, но и как более содержательный показатель надёжности.

Для оценки калибровки использовались две стандартные величины:

\begin{itemize}
\tightlist
\item
  \texttt{ECE} (Expected Calibration Error);
\item
  \texttt{Brier\ score}.
\end{itemize}

Калибровка анализировалась на preview-выборках по 5000 объектов. Сводные значения для validation и test приведены в
таблице~\ref{tab:calibration-metrics}.

\begin{center}
\small
\begin{longtable}{p{0.18\textwidth} r r r r r}
\caption{Калибровочные метрики на preview-выборках}
\label{tab:calibration-metrics}\\
\toprule
Split & ECE & Brier & Mean confidence & Accuracy@1 & \makecell[r]{Число\\объектов} \\
\midrule
\endfirsthead

\toprule
Split & ECE & Brier & Mean confidence & Accuracy@1 & \makecell[r]{Число\\объектов} \\
\midrule
\endhead

\bottomrule
\endlastfoot

validation & 0.02305 & 0.16440 & 0.75026 & 0.73560 & 5 000 \\
test & 0.01737 & 0.16120 & 0.75398 & 0.74800 & 5 000 \\

\end{longtable}
\footnotesize
Источник значений: \path{artifacts/results/recommendation/calibration_metrics_baseline.csv}.
\normalsize
\end{center}

Калибровочная кривая для test preview-выборки приведена на рисунке~\ref{fig:reliability-diagram-test}. Диагональная
линия соответствует идеальной калибровке, а точки показывают соотношение средней уверенности модели и фактической
точности внутри интервалов вероятности.

\begin{figure}[H]
\centering
\includegraphics[width=0.82\textwidth]{../../artifacts/results/recommendation/reliability_diagram_test_baseline.png}
\caption{Reliability diagram для baseline CatBoost на test preview-выборке}
\label{fig:reliability-diagram-test}
\end{figure}


Низкие значения \texttt{ECE} указывают на то, что вероятности модели в целом согласованы с фактической точностью на
рассматриваемых preview-выборках. Это не означает идеальную калибровку во всех областях вероятностного пространства, но
подтверждает, что confidence-сигнал можно использовать не только для сортировки кандидатов, но и как содержательный
показатель уверенности. Значения \texttt{Brier\ score} также поддерживают вывод о приемлемом качестве вероятностного
вывода модели.

\subsection{Место confidence-анализа в общей логике работы}

С архитектурной точки зрения confidence-слой является естественным продолжением recommendation layer. Сначала модель
формирует ранжированный список вероятных культур, затем confidence-анализ показывает, насколько устойчиво этот shortlist
выглядит с точки зрения вероятностного вывода. В результате система получает ещё один уровень интерпретации, не
связанный напрямую с семантическим объяснением кандидатов, но важный для практического использования.

Это делает confidence-анализ связующим звеном между recommendation-постановкой и explainability-слоем. Recommendation
layer отвечает на вопрос, какие кандидаты входят в shortlist. Confidence layer показывает, насколько надёжно этот
shortlist сформирован. Далее explainability-слой уже интерпретирует сами кандидаты и их взаимное положение внутри
списка. Такое разделение ролей делает общую архитектуру системы более прозрачной и позволяет не смешивать вероятностную
надёжность результата с его семантической интерпретацией.

\subsection{Выводы по главе}

В данной главе было показано, что рекомендательная система должна оцениваться не только по составу shortlist, но и по
степени уверенности модельного вывода. Для этого в работе введён confidence-слой, который разделяет случаи на группы
\texttt{high\_confidence}, \texttt{medium\_confidence} и \texttt{low\_confidence}. Результаты на test-выборке
показывают, что такое разделение действительно отражает различия в качестве рекомендаций: в высокоуверенных случаях
модель заметно точнее, а в низкоуверенных случаях сохраняет полезность прежде всего как shortlist-рекомендатель.

Дополнительный анализ калибровки показывает, что вероятности модели в целом согласованы с фактической точностью на
preview-выборках. Это усиливает практическую ценность confidence-сигнала и позволяет рассматривать его как
самостоятельный компонент системы поддержки принятия решений. Тем самым confidence-слой делает рекомендательный
результат более содержательным и подготавливает переход к следующему уровню анализа --- explainability через knowledge
graph.
